%old abstract
We introduce \bench, a new benchmark for evaluating AI agents for data-driven scientific discovery. 
Existing AI agents for science are usually evaluated on only a few tasks, leaving their practical utility in question. 
To bridge this gap, we collect 102 diverse, open-ended tasks in real-world data-driven discovery workflows.
We formulate each task as a challenging code generation problem and ask the agents to generate a complete program from scratch.
These programs are adapted from 44 peer-reviewed publications in four different domains, with code and data released under permissive licenses.
To further ensure their quality, the tasks are verified by seven domain experts and processed to prevent data contamination.
We evaluate five open-weight and proprietary LLMs on our benchmark with direct prompting, OpenDevin, and self-debug.
Even the best LLM-based agent can only solve $5.0\%$ of the tasks independently and $20.0\%$ of them with additional knowledge provided by domain experts.
These results highlight the need to further improve AI agents for data-driven scientific discovery, and \bench will facilitate the development process as a realistic and challenging testbed.

\hs{considering this flow: Recently there has been enormous interest in developing AI agents to automate scientific discovery. Concurrent work \cite{TheAIScientist} claims to have developed an LLM powered agent that can fully automate the research pipeline from generating hypothesis to running experiments to writing research papers, which has spurred heated discussions and received critical reviews of the true capabilities of the propose agent. In this paper, we will provide a better understanding on how far current language agents are from fully automating scientific discovery, by developing \bench, a new benchmark for evaluating AI agents to automate part of the workflow in data-driven scientific discovery. Specifically, we envision future agents to be \textit{scientist co-pilots} which can help domain scientists to write computer programs to automatically analyze data, build machine learning models, and visualize the findings. We examined 44 real peer-reviewed publications in four different scientific domains, from which we collected 102 diverse, open-ended tasks in data-driven discovery workflows and annotated the computer program to compete each task. To further ensure our task and program quality, they are verified by seven domain experts and processed to prevent data contamination \hs{need to rephrase and double check this sentence}. We evaluate five open-weight and proprietary LLMs on our benchmark with direct prompting, OpenDevin\hs{OpenHands?}, and self-debug \hs{this needs to be rephrased to convey that we evaluate both LLMs and their powered agentic frameworks.}.
Even the best LLM-based agent can only solve $5.0\%$ of the tasks independently and $20.0\%$ of them with additional knowledge provided by domain experts. These results highlight that \textit{current agents are far from being able to fully automate data-driven scientific discovery -- they even cannot automate part of the workflow yet, not to mention the whole research pipeline.} \bench will serve as a necessary (albeit not sufficient) benchmark to measure the progress towards language agents for automating scientific discovery.} %the need to further improve AI agents for data-driven scientific discovery, and \bench will facilitate the development process as a realistic and challenging testbed. 
\footnote{Code and data will be released online.}

%
\textit{What makes an AI agent practical\hs{useful?} for data-driven scientific discovery?}\hs{i think this question is not appropriate to motivate our work. The connection/flow between the second and the third paragraph could be strengthened. see related comments below:}
On one hand, \textbf{it shall play an assistive role to save scientists' time and efforts}.
While it is appealing to generate new research ideas or write papers with an agent,\hs{I think we could rephrase this sentence in a different tone: There are different goals for different AI agents for science; one desired goal is to generate new research ideas, conduct experiments and write papers autonomously, but } scientists would still have to carefully rule out potential issues, such as hallucination and inaccurate result interpretation, in the generated ideas and papers.
Instead, by \hs{here our goal is to} focusing on tedious tasks with easily verifiable outcomes, like data preparation, model development, data analysis, and information visualization (Figure \ref{fig:examples}), an agent can more effectively save the time and efforts of scientists. \hs{we should also add that it is an essential ability that an agent must have before they can fully autonomously conduct novel research in a scientific domain.}
On the other hand, \textbf{it shall work on diverse, open-ended tasks in real-world settings}.
These tasks can have multiple valid solutions and cover a wide range of datasets, including drug molecules, EEG time series, and geospatial maps (Figure \ref{fig:examples}).
Moreover, they come with only simple instructions but \textit{not} any detailed plans or code templates.
To perform well, an agent needs to reason about the complex workflow of each task, deal with heterogeneous data, and interact with the programming environment to adjust its code.
%

%
With these desiderata, we present \bench, a new benchmark consisting of 102 expert-validated tasks to evaluate LLM-based agents for data-driven scientific discovery. 
We construct the benchmark based on three principles:
First, \hs{I think this first point should be mentioned earlier when introducing the motivation and background, or merged to that part.} the tasks shall be feasible for an agent to carry out autonomously.
While there are agents developed to collaborate with scientists in wet-lab experiments \citep{coscientist,doi:10.1021/acscentsci.3c01087}, they are not yet capable of automating physical tasks and have to rely on humans. 
Thus, we focus on one paradigm for science, data-driven discovery, due to the digital nature of AI agents.
Second, to resemble real-world scientific discovery, \hs{this is our core novelty, and should be conveyed better.} we derive the tasks from open-source code and data of 44 peer-reviewed publications in four data-rich domains: Bioinformatics, Computational Chemistry, Geoinformatics, and Psychology \& Cognitive Neuroscience.
Since their code and data are available online, the tasks are also annotated with a procedure to prevent data contamination\hs{this is worth being talked about in a separate point} and other shortcuts \citep{kapoor2024aiagentsmatter} that may result in imprecise evaluation.
Finally, instead of an abstract workflow \citep{majumder2024discoverybenchdatadrivendiscoverylarge}, each task's output shall be easily verifiable and directly usable by a scientist without additional efforts to modify or implement. \hs{I feel this point about standalone Python program formulation should be talked about earlier. For example, right after the first sentence in this paragraph, you talk about the problem formulation and there you clarify the expected output is a python program and why. After that, you talk about the principles you hold in constructing this benchmark.}
To this end, we define the task outputs in our benchmark to be stand-alone Python programs. \hs{I think we should really stress on the expert efforts and multi-round validation to guarantee quality.}
%